% Preámbulo
\documentclass[stu, 12pt, letterpaper, floatsintext]{apa7}
\usepackage[utf8]{inputenc}
\usepackage{comment}
\usepackage{marvosym}
\usepackage{graphicx}
\usepackage{float}
\usepackage[normalem]{ulem}
\usepackage[spanish]{babel} 
\usepackage{lipsum}
\selectlanguage{spanish}
\usepackage{listings}
\usepackage{subcaption}
\usepackage{tabularx}
\usepackage{booktabs} % Para líneas profesionales (estilo APA)
\usepackage{mwe} % O simplemente asegúrate de usar \raisebox

\usepackage[style=apa, backend=biber]{biblatex}
\addbibresource{mibibliografia.bib} % Tu archivo de bibliografía
\useunder{\uline}{\ul}{}
\newcommand{\myparagraph}[1]{\paragraph{#1}\mbox{}\\}
\addto\captionsspanish{\renewcommand{\listtablename}{Índice de tablas}}
\addto\captionsspanish{\renewcommand{\lstlistlistingname}{Índice de códigos}}
\addto\captionsspanish{\renewcommand{\tablename}{Tabla}}
\addto\captionsspanish{\renewcommand{\lstlistingname}{Código}}

\providecommand{\inputpath}{./src/} % Carpeta sugerida para tus códigos
\DeclareGraphicsExtensions{.pdf,.jpeg,.png,.eps}
\graphicspath{ {./graphics/} {./img} }

% Colores estilo "Dark Mode" suave o editor moderno
\definecolor{codegreen}{rgb}{0,0.6,0}
\definecolor{codegray}{rgb}{0.5,0.5,0.5}
\definecolor{codepurple}{rgb}{0.58,0,0.82}
\definecolor{arduino_teal}{rgb}{0.0, 0.59, 0.62}
\definecolor{backcolour}{rgb}{0.95,0.95,0.92}

% ESTILO BASE (Con soporte de tildes incluido)
\lstdefinestyle{base_style}{
    backgroundcolor=\color{backcolour},   
    numberstyle=\tiny\color{codegray},
    basicstyle=\footnotesize\ttfamily,
    breaklines=true,                 
    numbers=left,                    
    numbersep=7pt,
    xleftmargin=10pt,                  
    showstringspaces=false,
    frame=single,
    rulecolor=\color{black!10},
    literate={á}{{\'a}}1 {é}{{\'e}}1 {í}{{\'i}}1 {ó}{{\'o}}1 {ú}{{\'u}}1
             {Á}{{\'A}}1 {É}{{\'E}}1 {Í}{{\'I}}1 {Ó}{{\'O}}1 {Ú}{{\'U}}1
             {ñ}{{\~n}}1 {Ñ}{{\~N}}1
             {α}{{\ensuremath{\alpha}}}1 {β}{{\ensuremath{\beta}}}1 
             {γ}{{\ensuremath{\gamma}}}1 {δ}{{\ensuremath{\delta}}}1
             {λ}{{\ensuremath{\lambda}}}1 {π}{{\ensuremath{\pi}}}1
             {Ω}{{\ensuremath{\Omega}}}1 {τ}{{\ensuremath{\tau}}}1
	% Soporte español
}

% ESTILOS ESPECÍFICOS
\lstdefinestyle{python_style}{
    style=base_style,
    language=Python,
    commentstyle=\color{codegreen},
    keywordstyle=\color{magenta},
    stringstyle=\color{codepurple}
}

\makeatletter
\newcommand{\mostrarautor}{\@author}
\makeatother

\makeatletter
\let\subparagraph\relax
\makeatother

\usepackage{titlesec}

% Configuración para alinear a la izquierda con numeración
\titleformat{\section}
	{\normalfont\large\bfseries}{\thesection.}{1em}{}
\titlespacing*{\section}{0pt}{12pt}{6pt}
    
\setcounter{secnumdepth}{3}